\section{Methods}\label{sec:methods}

\subsection{A construction card is an executable description of a release}

The unit of analysis is a person in a common row order. A population contains memberships, residential coordinates and, when available, parent and partner references. A rule set maps these inputs to layer-specific dyads and an eligibility mask. Let $X$ be the membership input, $R$ the resolved rules, and $s$ the seed. The build returns
\[
    B(X,R,s) = (I_R, A^{\rm household}, A^{\rm neighbour},
                      A^{\rm school}, A^{\rm work}, A^{\rm kin}),
\]
where each requested $A^\ell$ is a simple undirected adjacency matrix on the original row indices and $I_R$ marks eligible people. Excluded rows retain their index and have no edges. A user can therefore align two outputs without inferring an identity map from the edges. Age restrictions can still change the population over which descriptors are averaged.

The package ships \CardCount{} TOML cards. Table~\ref{tab:cards} lists their parent cards and changes. A card records values, source locations and an evidence status such as \emph{documented}, \emph{inferred} or \emph{assumed}. The loader merges base and child fields and constructs a frozen \texttt{RuleSet}. It rejects an unknown rule field and a cycle in the base chain. Cap values can name a size-dependent function. Provenance is read separately, so an unresolved source statement does not get silently interpreted by the build function. Unwritten fields use defaults. A provenance label records the card author's reading of a field. Table~\ref{tab:provenance-schema} shows the division between executable fields and metadata.
% Generated by src/make_tables.py.
\begin{table}[!htbp]
\centering\small
\caption{Bundled construction cards and their inheritance.}
\label{tab:cards}
\begin{tabular}{@{} >{\raggedright\arraybackslash}p{0.23\linewidth} >{\raggedright\arraybackslash}p{0.15\linewidth} >{\raggedright\arraybackslash}p{0.53\linewidth} @{}}
\toprule
\rowcolor{tableblue} \textbf{Card} & \textbf{Base} & \textbf{Implemented change} \\
\midrule
\texttt{dk-pairs} & \texttt{dk} & Changes the workplace draw to uniform pairs \\
\texttt{dk-twoway} & \texttt{dk} & Changes the workplace draw to random nominations and union \\
\texttt{dk} & \texttt{nl-2018} & Uses sites and the symmetric hard-cap draw \\
\texttt{nl-2018-classcap} & \texttt{nl-2018} & Adds a school nomination cap \\
\texttt{nl-2018} & --- & 2018 Dutch snapshot; reference for the separately documented releases \\
\texttt{nl-series} & \texttt{nl-2018} & Removes the neighbour radius and adds two kin categories \\
\texttt{se-code-nowindow} & \texttt{se-code} & Removes the graduation-cohort window \\
\texttt{se-code} & \texttt{se} & Substitutes the released link-function coefficients in work, property and school \\
\texttt{se} & --- & Eligibility, property, education and workplace rules specific to the Swedish approximation \\
\bottomrule
\end{tabular}
\end{table}

% Generated by src/make_tables.py.
\begin{table}[!htbp]
\centering\small
\caption{What goes into a construction card and what the loader does with it.}
\label{tab:provenance-schema}
\begin{tabular}{@{} >{\raggedright\arraybackslash}p{0.21\linewidth} >{\raggedright\arraybackslash}p{0.32\linewidth} >{\raggedright\arraybackslash}p{0.38\linewidth} @{}}
\toprule
\rowcolor{tableblue} \textbf{Entry} & \textbf{Author supplies} & \textbf{Loader behaviour} \\
\midrule
\texttt{[card]} & Name, release, sources, optional base and omissions & Resolves the base chain; the description combines source identifiers \\
\texttt{[rules]} & RuleSet fields and allowed literals & Merges fields, maps cap names to callables, rejects unknown fields \\
\texttt{[provenance]} & Field status and source location & Merges entries in base-to-child order; status is descriptive metadata \\
Unwritten fields & An explanation of relevant defaults & Falls back to the dataclass defaults; the loader does not infer evidence status \\
\bottomrule
\end{tabular}
\end{table}


\subsection{The bundled cards capture choices that matter across releases}

Table~\ref{tab:releases} shows the source rules by layer. The Dutch 2018 description groups workers by company. Companies with at most 101 employees produce cliques. In larger companies each worker nominates the 100 nearest colleagues by home location. The neighbourhood rule selects up to ten households within 50 metres and then links their members. Education groups are built from school, location, year and education type. The base card leaves school groups fully linked, in line with the out-degree distribution reported for the 2018 file. The alternative school cap represents another reading documented in applications \citep{vanderlaan2022,vanderlaan2023,bos2022}.
% Generated by src/make_tables.py.
\begin{table}[!htbp]
\centering\small
\caption{Published release definitions that motivate the cards. These are source descriptions.}
\label{tab:releases}
\begin{tabular}{@{} >{\raggedright\arraybackslash}p{0.15\linewidth} >{\raggedright\arraybackslash}p{0.18\linewidth} >{\raggedright\arraybackslash}p{0.18\linewidth} >{\raggedright\arraybackslash}p{0.18\linewidth} >{\raggedright\arraybackslash}p{0.21\linewidth} @{}}
\toprule
\rowcolor{tableblue} \textbf{Layer} & \textbf{Dutch 2018} & \textbf{Dutch yearly} & \textbf{Danish} & \textbf{Swedish} \\
\midrule
Household & Household identifier & Same address & Registered household & Co-resident relatives \\
Neighbour & Ten closest households within 50 m & Ten closest addresses plus a separate proximity relation & Ten closest households within 50 m & Same property, size-dependent draws \\
School & School, location, type, year & Release-specific education levels & Grade and programme, or entry cohort & Level, municipality and graduation year \\
Work & Company and nearest homes & Jobs pooled over months; source conflict over cap scope & Workplace, large-group sampling & Company or workplace, size-dependent draws \\
Kin & Legal parent links & Adds partners and extended relations & Legal parent links & Co-residence and age differences \\
\bottomrule
\end{tabular}
\end{table}


The yearly Dutch documentation treats addresses as the residential unit and separates nearest-address neighbours from a second proximity relation \citep{cbs2023,das2023,catbur}. It also documents address-based co-residence, changes in education grouping and an expanded family relation \citep{cathui,catkla,catfam}. The colleague codebook pools jobs over the year, with a conflict over the scope of the final cap when compared with the technical appendix to a later CBS application \citep{catcol,vanklingeren2024}. The card removes the distance limit using households, and adds partner and in-law categories. It omits the second neighbour relation, monthly jobs, full address structure and several kin categories. It inherits household cliques in place of the release-specific changes to the household layer.

The Danish card groups workers by workplace site and inherits a household-neighbour rule with a radius. The source description establishes workplace sampling and symmetrisation, while the procedure used for large workplaces remains uncertain \citep{cremers2025}. Three executable alternatives clarify the range of cap definitions. A hard cap restrains every realised vertex degree. A random-pairs rule fixes the number of distinct pairs. A nomination rule lets each person draw up to the nomination limit and includes the union. The latter can give each person more neighbours than they named. The card uses the same simplified school keys as the Dutch baseline. It omits Danish programme and tertiary cohort definitions, fictional workplaces and the provision for flatmates.

The Swedish card applies a minimum age of 15, splits companies of at least 200 workers into workplaces, uses property neighbours and groups completed education by graduation cohort. The thesis and released code motivate two versions of the size-dependent nomination count. The code variant changes the coefficients in work, property and school. A further variant removes the assumed 19-year cohort look-back. A cohort size limit of 10,000 remains. The source school definition includes a field-of-study split and the work construction includes special handling of workplaces. The cards omit both details. The household operator still forms household cliques, whereas the source household layer is restricted to co-resident relatives. The Swedish source distinguishes close and extended family. The package uses one family layer and filters the available parent and partner links using age and residence proxies \citep{panayiotou2025,wohlert2024,secode2026}.

\subsection{Membership tables and projection operators use stable person indices}

\texttt{world\_from\_tables} accepts arrays in row order and remaps non-negative group identifiers to consecutive integers where needed. A household identifier is required for each person. The caller supplies projected coordinates in metres for spatial operations, and household coordinates in sorted identifier order if supplied separately. Optional grouping keys use negative sentinels for missing membership. A missing site vector defaults to employer membership. If both coordinate arrays are omitted, the adapter uses coincident household locations. Omitted age defaults to adulthood, and missing school level for a recorded school defaults to the primary level. In an applied analysis these defaults require a deliberate input decision, since they can change eligibility or tie construction. Table~\ref{tab:inputs} gives information each family of rules requires.
% Generated by src/make_tables.py.
\begin{table}[!htbp]
\centering\small
\caption{Membership input contract for \texttt{world\_from\_tables}.}
\label{tab:inputs}
\begin{tabular}{@{} >{\raggedright\arraybackslash}p{0.25\linewidth} >{\raggedright\arraybackslash}p{0.32\linewidth} >{\raggedright\arraybackslash}p{0.34\linewidth} @{}}
\toprule
\rowcolor{tableblue} \textbf{Input} & \textbf{Required to express} & \textbf{Important convention} \\
\midrule
Household IDs & All projections; one person per row & Use a separate ID for every single-person household \\
Home or household coordinates & Nearest-home work and neighbourhoods & Use projected metres; household rows follow sorted unique household IDs \\
Employer and site IDs & Company and establishment grouping & Supply globally distinct site IDs; missing sites can collapse to employer \\
School IDs, level, year, class & Current enrolment groups & Unknown category keys can exclude a member \\
Property IDs & Property-based neighbours & Use the plot or property definition for the target release \\
Municipality, attainment, graduation & Completed education groups & Graduation is recorded as years before the snapshot \\
Age and kin references & Eligibility and relation composition & Kin references are zero-based input row indices \\
\bottomrule
\end{tabular}
\end{table}


The constructor processes layers in a fixed order. It derives one seeded random stream per layer from the run seed and layer index. A change in one layer therefore does not consume another layer's draws. Within a capped group it makes nominations by distance or random selection, or constructs a constrained set of unordered pairs. For a fixed nomination cap $k$, a group of size at most $k+1$ becomes a clique by default. The sequential cap chooses eligible partners in random order and removes a vertex from the eligible pool when it reaches its limit. The random-pair operator samples distinct pairs. The household-neighbour operator searches at household resolution and expands each selected household pair to the Cartesian product of its members. For nearest searches the implementation adds a small seeded coordinate perturbation before querying a spatial tree. A point near a radius boundary can therefore be treated differently from a distance test on the original coordinates.

Kin construction composes parent, partner and late-parent links with sparse operations. Selected parent, sibling, grandparent, aunt or uncle, cousin and co-parent relations can be combined. The yearly Dutch implementation can also add partners and in-laws. The operator uses intermediary relationships to derive kin before the final eligibility filter. The output for each layer combines repeated nominations and drops self-loops. In particular, if either person names the other, the pair becomes one edge in the layer. Directed nomination records are discarded. Table~\ref{tab:operators} describes the bounds that follow from these choices. Full linking emits a quadratic number of pairs in a group. The uniform-nomination operator emits at most $nk$ pairs from a capped group but draws random keys over all $n$ candidates per member in batches. That step still evaluates a quadratic number of keys. The sequential cap may also scan the full membership for each processed vertex. Uniform pairs instead sample ranks from the upper triangle and decode them to endpoints. These implementation differences matter even when two operators have similar mean degree.
% Generated by src/make_tables.py.
\begin{table}[!htbp]
\centering\small
\caption{Local projection operators and bounds on the simple undirected output.}
\label{tab:operators}
\begin{tabular}{@{} >{\raggedright\arraybackslash}p{0.21\linewidth} >{\raggedright\arraybackslash}p{0.37\linewidth} >{\raggedright\arraybackslash}p{0.33\linewidth} @{}}
\toprule
\rowcolor{tableblue} \textbf{Operator} & \textbf{Choice} & \textbf{Consequence} \\
\midrule
Full linking & All unordered co-member pairs & $d=n-1$ in one group \\
Nearest nomination & Each person names $k$ nearest co-members & Union can exceed $k$ degree \\
Random nomination & Each person samples $k$ distinct other people & Union has mean $2k-k^2/(n-1)$ above the clique threshold \\
Sequential hard cap & Process vertices in a random order and draw among eligible others & No vertex exceeds the cap in a sampled group \\
Random pairs & Sample $\lfloor nk/2\rfloor$ distinct unordered pairs & Mean is near the cap; individuals may exceed it \\
Household neighbours & Nearest household search then expand linked pairs & A household nomination can produce many person ties \\
Kin composition & Sparse products of directed parent incidence and related matrices & Outputs the union over selected categories \\
\bottomrule
\end{tabular}
\end{table}


\subsection{Sparse adjacency supports both union and multiplex descriptors}

The core uses NumPy arrays and SciPy compressed sparse row matrices \citep{harris2020,virtanen2020}. The matrices store two entries per undirected edge, with duplicate proposals reduced to binary adjacency. The sum of layer matrices is the multiplexity matrix $M$. The union adjacency $U$ has an entry wherever $M$ is positive. This distinction prevents degree from counting a colleague who is also a neighbour twice. Clustering uses $U$ and is conditional on an ego having at least two distinct neighbours. Triangle counts use batches of rows of $(U^2)\circ U$. For a symmetric binary matrix, each closed ego-centred alter pair contributes twice to the corresponding row sum. Inactive rows are masked in the output. Sparse storage avoids allocating a full adjacency array. The intermediate products can still contain many two-step paths when groups are dense, so equal edge counts need not imply equal descriptor times. Table~\ref{tab:descriptors} summarises the network and affiliation views.
% Generated by src/make_tables.py.
\begin{table}[!htbp]
\centering\small
\caption{Public outputs and the representation each interprets.}
\label{tab:descriptors}
\begin{tabular}{@{} >{\raggedright\arraybackslash}p{0.26\linewidth} >{\raggedright\arraybackslash}p{0.3\linewidth} >{\raggedright\arraybackslash}p{0.35\linewidth} @{}}
\toprule
\rowcolor{tableblue} \textbf{Output} & \textbf{Computed on} & \textbf{Interpretation} \\
\midrule
Layer or union degree & Per-layer binary adjacency or its union & Number of distinct neighbours in the chosen view \\
Local clustering & Binary union & Fraction of alter pairs adjacent, among egos with at least two alters \\
Excess closure & Layer matrices and multiplexity & Normalised gain over pure-layer closure under a named convention \\
Multiplex share & Multiplexity counts & Share of an ego’s neighbours in multiple layers \\
Participation & Layer degree vector & Evenness of degree across layers \\
Affiliation summaries & Membership arrays and coordinates & Group sizes, cross-context membership and distances before projection \\
\bottomrule
\end{tabular}
\end{table}


Excess closure has a further choice. For an ego $u$ with summed layer degree $K_u$ and multiplicities $m_{uv}$, the number of pairs of incident ties to distinct alters is
\[
 P_u = \binom{K_u}{2}-\sum_v \binom{m_{uv}}{2}.
\]
Let $T_{\rm pure}$ be the sum of within-layer triangles at $u$. With $c_{\rm pure}=T_{\rm pure}/P_u$ and $c_{\rm act}=T/P_u$, the excess is $(c_{\rm act}-c_{\rm pure})/(1-c_{\rm pure})$. The \texttt{paper} numerator counts one per adjacent pair of alters. The \texttt{mlnlib} numerator weights that pair by $m_{uv}m_{vw}m_{wu}$. The \texttt{tiepair} numerator weights it by $m_{uv}m_{uw}$ and requires a tie between the two alters. A zero denominator gives a missing value. The stored example in Section~\ref{res:closure} distinguishes the three definitions.

When both populations can be built under both cards, the package also supplies an arithmetic decomposition. Write $m_{ij}$ for a summary of population $i$ under card $j$. The naive difference is the difference between $m_{22}$ and $m_{11}$. The construction component is the mean of $m_{12}-m_{11}$ and $m_{22}-m_{21}$. The population component is the mean of $m_{21}-m_{11}$ and $m_{22}-m_{12}$. Their sum equals the diagonal difference. This is a symmetric standardisation of the comparison \citep{kitagawa1955,dasgupta1993,shorrocks2013}. A large interaction means that the rule contrast itself depends on the population. The affiliation module offers separate descriptions of group size, context overlap and residential distance before any edge construction.

\subsection{Analytic expectations use separate inputs from the graph builder}

Full linking yields degree $n-1$ in a group of size $n$ and clustering one for $n\geq3$. Suppose instead that members independently nominate $k<n-1$ distinct peers uniformly. With $q=k/(n-1)$, a particular pair is absent with probability $(1-q)^2$, so the exact mean union degree is $2k-k^2/(n-1)$. This reasoning allows dependence among pair indicators. The package uses $1-(1-q)^2$ as a clustering approximation. It is the marginal pair probability, while local clustering also conditions on edges adjacent to the ego. Fixed-choice analyses have used related nomination-and-union models \citep{momeni2017}.

The nearest-home and sequential hard-cap rules use stored tabulations. For spatial neighbours the model uses the expected number $\mu$ of other households within the radius. In a homogeneous Poisson setting, the probability of no other household in that radius is $e^{-\mu}$. Every household with at least one candidate inside the radius nominates one or more households when its cap is positive. An isolated household in this setting cannot receive a nomination either. The package interpolates the saved linked-household counts. Archived neighbourhood clustering is available for display and interpretation, while the public theory functions for the neighbourhood return linked-household degree and isolation. Table~\ref{tab:theory} distinguishes exact expectations, proxies and interpolation.
% Generated by src/make_tables.py.
\begin{table}[!htbp]
\centering\small
\caption{How the package obtains edge-free group and household summaries.}
\label{tab:theory}
\begin{tabular}{@{} >{\raggedright\arraybackslash}p{0.22\linewidth} >{\raggedright\arraybackslash}p{0.33\linewidth} >{\raggedright\arraybackslash}p{0.37\linewidth} @{}}
\toprule
\rowcolor{tableblue} \textbf{Rule} & \textbf{Degree or isolation} & \textbf{Within-layer clustering} \\
\midrule
Clique & $n-1$ for a group of size $n$ & $1$ when $n\geq3$ \\
Random nominations & Exact expectation $2k-k^2/(n-1)$ & Uses marginal tie probability as an approximation \\
Random pairs & Approximately $k$; exact mean $2\lfloor nk/2\rfloor/n$ & Uses pair density as an approximation \\
Sequential cap & Interpolates degree from the cap-100 grid & Interpolates cap-100 clustering; density proxy for other caps \\
Spatial nearest & Interpolation over the archived spatial grids & Separate curves by coordinate model \\
Household radius & Interpolated linked-household count, $\exp(-\mu)$ isolation & Archived grid under homogeneous settlement \\
\bottomrule
\end{tabular}
\end{table}


Layer mean degree from a size distribution weights each group by its member count and includes zero contributions for those outside groups. The clustering aggregator weights the groups of at least three members. For nearest-home groups the archived coordinate patterns are uniform on a square or Gaussian. The interpolator works in log excess group size relative to the nomination cap, uses the nearest available tabulated cap for a new one, and holds the last value beyond the grid. These choices give a low-input approximation. The link with geometric nearest-neighbour reciprocity is central because the number of mutual nominations determines how many directed choices collapse to a single undirected edge \citep{schilling1986}.

\subsection{Validation distinguishes code checks, demonstrations and pre-specified predictions}

Table~\ref{tab:questions} states the question attached to each evaluation. The release commits fix the operators, built-in scenarios and existing regression assertions before this manuscript. The regression suite tests software behaviour and has no registered empirical outcome. Parity against a separate triangle computation is also a code check. The available aggregate predictions were pre-specified in the companion study before its analyses, conditional on already published targets. These checks use public aggregates. The Danish alternative calculations followed inspection of the failed prediction. The later Danish sensitivity and theory amendments were written after inspection of earlier outcomes in the same research programme and used distinct evaluation runs where specified. This article uses the release code plus their archived evidence, without refitting or changing the original verdicts. Supplement~\ref{sec:provenance} gives the order, dates and deviations.
% Generated by src/make_tables.py.
\begin{table}[!htbp]
\centering\small
\caption{Validation questions and the strongest evidence available for each.}
\label{tab:questions}
\begin{tabular}{@{} >{\raggedright\arraybackslash}p{0.25\linewidth} >{\raggedright\arraybackslash}p{0.35\linewidth} >{\raggedright\arraybackslash}p{0.31\linewidth} @{}}
\toprule
\rowcolor{tableblue} \textbf{Question} & \textbf{Evidence} & \textbf{Scope} \\
\midrule
Q1. Do the operators implement their stated choices? & Release tests, hand-worked fixtures and triangle parity & Code behaviour and comparison implementations \\
Q2. Can users compare cards on one input? & Existing example evaluated for all four reference cards & Controlled input, card-dependent active mask \\
Q3. Can layer summaries be estimated without edges? & Exact identities and archived simulation grids & Group and household models, conditional on assumptions \\
Q4. Can ambiguities be diagnosed? & Saved companion predictions and a two-layer closure example & One aggregate comparison and a definitional counterexample \\
Q5. What does local execution cost? & Stage wall times at the two built-in scales & Descriptive software measurements on one system \\
\bottomrule
\end{tabular}
\end{table}


The runtime evaluation is descriptive. We timed the README example at its built-in \texttt{small} setting and the corresponding existing \texttt{main} preset, with membership and wiring seed \ExampleSeed{} and \RepeatCount{} repetitions. The cards were run in the README order. Each repetition generated one population per scale and reused it for all four builds. Wall-clock time was recorded separately for generation, graph building and the default \texttt{position\_descriptors} call. Card loading, plotting and summarising results were outside the timed stages. The measurements used an allocation of \CpuCount{} cores with the numerical-library thread limit set to \ThreadCount{}. The median and range summarise the fixed set of executions. There was no timing-based stopping decision or package tuning. The two presets alter geography, employer limits and the number of households together.
\FloatBarrier
